\section{Variáveis e operadores}

\begin{frame}[fragile,t]{1. Variáveis}
O que são? Variáveis são nomes associados a valores armazenados durante a execução do programa. Em Python, não é necessário declarar previamente o tipo da variável.\par
\begin{columns}[T,onlytextwidth]
\begin{column}{0.46\textwidth}
\EMASCcodelabel{Estrutura do comando}
\begin{lstlisting}
nome_da_variavel = valor
\end{lstlisting}
\end{column}
\begin{column}{0.51\textwidth}
\begin{visibleenv}<2->
\EMASCcodelabel{Exemplo prático}
\begin{lstlisting}
nome = "Carlos"
idade = 20
altura = 1.75

print(nome)
print(idade)
print(altura)
\end{lstlisting}
\end{visibleenv}
\end{column}
\end{columns}
\begin{visibleenv}<3->
\begin{EMASCnote}
Neste exemplo, nome guarda um texto, idade guarda um inteiro e altura guarda um número de ponto flutuante. O sinal \texttt{=} atribui um valor à variável.
\end{EMASCnote}
\end{visibleenv}
\end{frame}

\begin{frame}[t]{2. Tipos de variáveis}
O que são? O tipo indica a categoria do dado armazenado. Os tipos básicos mais comuns são:\par
\begin{tabularx}{\linewidth}{@{}lXl@{}}
\toprule
\textbf{Tipo} & \textbf{Descrição} & \textbf{Exemplo} \\
\midrule
\texttt{int} & Número inteiro & \texttt{idade = 18} \\
\texttt{float} & Número de ponto flutuante & \texttt{altura = 1.75} \\
\texttt{str} & Texto (string) & \texttt{nome = "Ana"} \\
\texttt{bool} & Valor lógico: verdadeiro ou falso & \texttt{aprovado = True} \\
\bottomrule
\end{tabularx}
\par\smallskip
Estrutura: atribua um valor e utilize \texttt{type()} para descobrir seu tipo.\par
\end{frame}

\begin{frame}[fragile,t]{2. Tipos de variáveis / Exemplo}
\begin{columns}[T,onlytextwidth]
\begin{column}{0.46\textwidth}
\EMASCcodelabel{Estrutura do comando}
\begin{lstlisting}
variavel = valor
type(variavel)
\end{lstlisting}
\end{column}
\begin{column}{0.51\textwidth}
\begin{visibleenv}<2->
\EMASCcodelabel{Exemplo prático}
\begin{lstlisting}
idade = 18
altura = 1.70
nome = "João"
aprovado = True

print(type(idade))
print(type(altura))
print(type(nome))
print(type(aprovado))
\end{lstlisting}
\end{visibleenv}
\end{column}
\end{columns}
\begin{visibleenv}<3->
\begin{EMASCnote}
A função \texttt{type()} retorna o tipo do objeto. Os resultados serão, respectivamente, \texttt{int}, \texttt{float}, \texttt{str} e \texttt{bool}.
\end{EMASCnote}
\end{visibleenv}
\end{frame}

\begin{frame}[fragile,t]{3. INPUT e PRINT}
O que é? \texttt{input()} recebe uma informação digitada pelo usuário; \texttt{print()} exibe informações na tela.\par
\begin{columns}[T,onlytextwidth]
\begin{column}{0.40\textwidth}
\EMASCcodelabel{Estrutura do comando}
\begin{lstlisting}
variavel = input("Mensagem: ")
print(variavel)
\end{lstlisting}
\end{column}
\begin{column}{0.56\textwidth}
\begin{visibleenv}<2->
\EMASCcodelabel{Exemplo prático}
\begin{lstlisting}
nome = input("Digite seu nome: ")
idade = int(input("Digite sua idade: "))

print("Seu nome é:", nome)
print("Sua idade é:", idade)
\end{lstlisting}
\end{visibleenv}
\end{column}
\end{columns}
\begin{visibleenv}<3->
\begin{EMASCnote}
Entenda o exemplo: O \texttt{input()} retorna texto (\texttt{str}). Por isso, \texttt{int()} converte a idade para inteiro. Para números de ponto flutuante, utiliza-se \texttt{float()}.
\end{EMASCnote}
\end{visibleenv}
\end{frame}

\begin{frame}[t]{4.1 Operadores aritméticos}
O que são? Operadores realizam cálculos, comparações e combinações lógicas.\par
\begin{tabularx}{\linewidth}{@{}lX@{}}
\toprule
\textbf{Operador} & \textbf{Operação} \\
\midrule
\texttt{+} & Adição \\
\texttt{-} & Subtração \\
\texttt{*} & Multiplicação \\
\texttt{/} & Divisão \\
\texttt{//} & Divisão inteira \\
\texttt{\%} & Resto da divisão \\
\texttt{**} & Potenciação \\
\bottomrule
\end{tabularx}
\end{frame}

\begin{frame}[fragile,t]{4.1 Operadores aritméticos / Exemplo}
\begin{columns}[T,onlytextwidth]
\begin{column}{0.48\textwidth}
\EMASCcodelabel{Estrutura do comando}
\begin{lstlisting}
resultado = valor1 operador valor2
\end{lstlisting}
\end{column}
\begin{column}{0.49\textwidth}
\begin{visibleenv}<2->
\EMASCcodelabel{Exemplo prático}
\begin{lstlisting}
numero1 = 10
numero2 = 5

print(numero1 + numero2) # 15
print(numero1 - numero2) # 5
print(numero1 * numero2) # 50
print(numero1 / numero2) # 2.0
print(numero1 % numero2) # 0
print(numero1 ** numero2) # 100000
\end{lstlisting}
\end{visibleenv}
\end{column}
\end{columns}
\end{frame}

\begin{frame}[fragile,t]{4.2 Operadores relacionais}
Comparam valores e produzem um resultado booleano (\texttt{True} ou \texttt{False}): \texttt{==} igual; \texttt{!=} diferente; \texttt{\textgreater{}} maior; \texttt{\textless{}} menor; \texttt{\textgreater{}=} maior ou igual; \texttt{\textless{}=} menor ou igual.\par
\begin{visibleenv}<2->
\EMASCcodelabel{Exemplo prático}
\begin{lstlisting}
idade = 18

print(idade >= 18)    # True
print(idade == 18)    # True
print(idade != 20)    # True
\end{lstlisting}
\end{visibleenv}
\end{frame}

\begin{frame}[fragile,t]{4.3 Operadores lógicos}
\texttt{and} (E) exige que as condições sejam verdadeiras; \texttt{or} (OU) exige pelo menos uma condição verdadeira; \texttt{not} (NÃO) inverte o valor lógico.\par
\begin{visibleenv}<2->
\EMASCcodelabel{Exemplo prático}
\begin{lstlisting}
idade = 20
tem_carteira = True

print(idade >= 18 and tem_carteira)     # True
\end{lstlisting}
\end{visibleenv}
\end{frame}
